
\documentclass[12pt,a4paper]{article}
\usepackage[utf8]{inputenc}
\usepackage{tabularx}
\usepackage[vietnamese]{babel}
\usepackage{geometry}
\geometry{margin=2.5cm}
\setlength{\parindent}{0pt}

\title{50 Câu Hỏi Trắc Nghiệm Về CSS Selectors}
\date{}

\begin{document}
\maketitle

\section*{Câu hỏi}

\textbf{Câu 1:} CSS selector nào chọn tất cả các phần tử?

\begin{tabularx}{\textwidth}{X X X X}
    A. \texttt{\#*} & B. \texttt{all} & C. \texttt{*} & D. \texttt{.}
\end{tabularx}

\textbf{Câu 2:} Selector nào chọn phần tử có id="header"?

\begin{tabular}{llll}
    A. .header & B. \#header & C. header & D. *header
\end{tabular}

\textbf{Câu 3:} Selector nào chọn tất cả phần tử có class="box"?
\begin{tabular}{llll}
A. \#box &
B. .box &
C. box &
D. *box
\end{tabular}

\textbf{Câu 4:} Selector nào chọn tất cả thẻ p?
\begin{tabular}{llll}
    A. .p &
    B. \#p &
    C. p &
    D. *p
\end{tabular}

\textbf{Câu 5:} Selector nào chọn tất cả p bên trong div?

\begin{tabular}{llll}
    A. div + p &
    B. div > p &
    C. div p &
    D. div ~ p 
\end{tabular}

\textbf{Câu 6:} Selector nào chỉ chọn con trực tiếp p của div?
\begin{tabular}{llll}
    A. div p &
    B. div > p &
    C. div + p &
    D. div ~ p \\
\end{tabular}

\textbf{Câu 7:} Selector nào chọn p đứng ngay sau div?
A. div > p \\
B. div p \\
C. div + p \\
D. div ~ p \\

\textbf{Câu 8:} Selector nào chọn tất cả p sau div cùng cấp?
A. div + p \\
B. div > p \\
C. div p \\
D. div ~ p \\

\textbf{Câu 9:} Selector nào chọn input có type="text"?
A. input(text) \\
B. input[type=text] \\
C. input=text \\
D. input\#text \\

\textbf{Câu 10:} Selector nào chọn phần tử có thuộc tính disabled?
A. [disabled] \\
B. .disabled \\
C. \#disabled \\
D. disabled \\

\textbf{Câu 11:} a:hover là loại selector nào?
A. Class \\
B. Pseudo-element \\
C. Pseudo-class \\
D. Attribute \\

\textbf{Câu 12:} Pseudo-element chọn ký tự đầu tiên:
A. :first \\
B. ::first-letter \\
C. :first-letter \\
D. ::first-line \\

\textbf{Câu 13:} Selector nào chọn phần tử đầu tiên?
A. :first \\
B. :first-child \\
C. :nth-child(1) \\
D. B và C \\

\textbf{Câu 14:} Selector chọn phần tử cuối:
A. :last \\
B. :last-child \\
C. :end \\
D. :final \\

\textbf{Câu 15:} p.intro nghĩa là:
A. p trong intro \\
B. p có id intro \\
C. p có class intro \\
D. p và intro \\

\textbf{Câu 16:} .box.large nghĩa là:
A. box trong large \\
B. large trong box \\
C. phần tử có cả box và large \\
D. box hoặc large \\

\textbf{Câu 17:} Class bắt đầu bằng col-:
A. \texttt{[class\^="col-"]} 
B. \texttt{[class*="col-"]} 
C. \texttt{[class\$="col-"]} 
D. \texttt{[class="col-"]} 

\textbf{Câu 18:} [class\$="end"] nghĩa là:
A. bắt đầu bằng end \\
B. chứa end \\
C. kết thúc bằng end \\
D. bằng end \\

\textbf{Câu 19:} :nth-child(2n) chọn:
A. phần tử lẻ \\
B. phần tử chẵn \\
C. phần tử thứ 2 \\
D. phần tử cuối \\

\textbf{Câu 20:} Input đang focus:
A. input:hover \\
B. input:active \\
C. input:focus \\
D. input:selected \\

\textbf{Câu 21:} Selector nhóm:
A. h1 h2 h3 \\
B. h1 + h2 + h3 \\
C. h1, h2, h3 \\
D. h1 > h2 > h3 \\

\textbf{Câu 22:} Chọn tất cả ngoại trừ p:
A. :not(p) \\
B. !p \\
C. except(p) \\
D. p:not \\

\textbf{Câu 23:} Tạo nội dung trước phần tử:
A. :before \\
B. ::before \\
C. :start \\
D. ::start \\

\textbf{Câu :} Độ ưu tiên cao nhất:
A. Element \\
B. Class \\
C. ID \\
D. Universal \\

\textbf{Câu :} Universal selector:
A. . \\
B. \# \\
C. * \\
D. \% \\

\textbf{Câu :} Chọn a có href:
A. a.href \\
B. a[href] \\
C. a\#href \\
D. a=href \\

\textbf{Câu :} div > * nghĩa là:
A. mọi div \\
B. mọi con trực tiếp của div \\
C. mọi cháu \\
D. div kế tiếp \\

\textbf{Câu :} :checked dùng cho:
A. input text \\
B. checkbox/radio \\
C. button \\
D. textarea \\

\textbf{Câu :} p ~ span chọn:
A. span trong p \\
B. span trước p \\
C. span cùng cấp sau p \\
D. span con p \\

\textbf{Câu :} Selector chọn phần tử rỗng:
A. :blank \\
B. :empty \\
C. :null \\
D. :void \\

\textbf{Câu :} :disabled chọn:
A. phần tử ẩn \\
B. phần tử bị khóa \\
C. hover \\
D. active \\

\textbf{Câu :} :enabled chọn:
A. phần tử có thể dùng \\
B. bị khóa \\
C. ẩn \\
D. lỗi \\

\textbf{Câu :} li:first-child nghĩa là:
A. li đầu của cha \\
B. li mọi nơi \\
C. li class first \\
D. li cuối \\

\textbf{Câu :} li:nth-child(3) chọn:
A. li thứ 2 \\
B. li thứ 3 \\
C. li thứ 4 \\
D. li cuối \\

\textbf{Câu :} div[class] chọn:
A. div có class \\
B. div không class \\
C. mọi div \\
D. div có id \\

\textbf{Câu :} [type="password"] là:
A. class selector \\
B. id selector \\
C. attribute selector \\
D. element selector \\

\textbf{Câu :} section article p là:
A. p con trực tiếp article \\
B. p trong article trong section \\
C. p sau article \\
D. p trước article \\

\textbf{Câu :} h1 + p chọn:
A. mọi p \\
B. p trong h1 \\
C. p ngay sau h1 \\
D. p trước h1 \\

\textbf{Câu :} input:required chọn:
A. input trống \\
B. input bắt buộc \\
C. input ẩn \\
D. input khóa \\

\textbf{Câu :} :root đại diện:
A. body \\
B. html \\
C. head \\
D. document \\

\textbf{Câu :} ::after là:
A. pseudo-class \\
B. attribute \\
C. pseudo-element \\
D. element \\

\textbf{Câu :} img[alt] chọn:
A. ảnh có alt \\
B. ảnh không alt \\
C. mọi ảnh \\
D. ảnh lỗi \\

\textbf{Câu :} :nth-of-type(2) chọn:
A. phần tử thứ 2 bất kỳ \\
B. phần tử thứ 2 cùng loại \\
C. phần tử cuối \\
D. phần tử đầu \\

\textbf{Câu :} button:hover khi:
A. click \\
B. giữ chuột \\
C. di chuột vào \\
D. rời chuột \\

\textbf{Câu :} .menu li a chọn:
A. mọi a \\
B. a trong li trong menu \\
C. li trong menu \\
D. menu trong a \\

\textbf{Câu :} * có độ ưu tiên:
A. cao nhất \\
B. trung bình \\
C. thấp nhất \\
D. bằng class \\

\textbf{Câu :} Kết hợp ID và class:
A. \#box .item \\
B. \#box.item \\
C. .box\#item \\
D. item\#box \\

\textbf{Câu :} input:valid chọn:
A. input sai \\
B. input đúng định dạng \\
C. input rỗng \\
D. input khóa \\

\end{document}
